% ============================================================
\section{The Balanced Dimension and $K_4$}
\label{sec:balanced-dimension}
% ============================================================

The holographic violation $\cH(D) = \binom{D}{2} - (D{-}1)!$ measures the
mismatch between pairwise couplings and orderings at an incidence of~$D$
primitives. $\cH(D)=0$ selects the balanced dimensions where these two counts
agree. This is a Diophantine equation with exactly two solutions: $D=2$
(degenerate) and $D=4$ (the unique non-degenerate balanced dimension). The
balanced complex~$K_4$ carries a six-slot edge carrier that decomposes as
$1{+}2{+}3$ under the canonical~$S_4$ action. The representation theory
of~$S_4$ is standard \cite{Serre1977,FultonHarris1991}; the intrinsic
decomposition uses only the character table and the edge carrier.

\subsection{The Completeness Equation}
\label{subsec:completeness-equation}
% ------------------------------------------------------------

\begin{definition}[Dimension]\label{def:dimension}\lean{dimensionOfEvent}
The \emph{dimension} $D$ of an incidence is the count of primitives meeting there.
\end{definition}

\begin{definition}[Primitive Normalization]\label{def:fundamental-scale}\lean{fundamentalScale}
A single primitive fixes the unit normalization. All later counting scales are measured relative to that unit.
\end{definition}

At an incidence of dimension $D$, structure is counted two ways:

\textbf{Pair count:} Pairwise couplings give $\binom{D}{2}$ unordered pairs.

\textbf{Order count:} With one primitive distinguished, the remaining $D-1$ can be ordered in $(D-1)!$ ways.

\begin{theorem}[Completeness]\label{thm:completeness}\lean{completeness_from_closure}\lean{stable_D2}\lean{stable_D4}
If an incidence of dimension $D$ is counted so that its pair count is $\binom{D}{2}$ and its order count is $(D-1)!$, then Closure Balance forces:
\begin{equation}
\boxed{\binom{D}{2} = (D-1)!}
\label{eq:completeness}
\end{equation}
\end{theorem}

\begin{proof}
Closure Balance identifies the pair count and order count for the incidence. Under the counting model used here, those counts are $\binom{D}{2}$ and $(D-1)!$, so the displayed equation follows immediately.
\end{proof}

\begin{definition}[Closure Residual]\label{def:holographic-operator}\lean{holographicOperator}
$\cH(D) := \binom{D}{2} - (D-1)!$
\end{definition}

\begin{definition}[Counting realization of the closure defect]
\label{def:counting-realization-closure-defect}\lean{CoherenceCalculus.countingRealizationClosureDefect}
Let \(I\) be a pure-dimensional incidence of dimension \(D\) in the sense of
Definition~\ref{def:dimension}. Its \emph{counting realization} is
\[
\lambda(I):=\cH(D).
\]
This is the first scalar counting realization of the abstract closure defect
\(\Lambda(I)\) from Definition~\ref{def:closure-defect}.
\end{definition}

\begin{proposition}[Pure-dimensional counting realization]
\label{prop:pure-dimensional-counting-realization}\lean{CoherenceCalculus.pure_dimensional_counting_realization}
If \(I\) is a pure-dimensional incidence of dimension \(D\), then
\[
\lambda(I)=\cH(D)=\binom{D}{2}-(D-1)!.
\]
In particular, the counting realization vanishes exactly on the
closure-balanced dimensions and is nonzero exactly on the counting residual
dimensions.
\end{proposition}

\begin{proof}
The first identity is Definition~\ref{def:counting-realization-closure-defect}.
The second statement is immediate from the definition of \(\cH(D)\).
\end{proof}

\begin{proposition}[Counting closure realizes closure balance on the pure-dimensional surface]
\label{prop:counting-closure-realizes-terminality}\lean{CoherenceCalculus.counting_zero_implies_balanced, CoherenceCalculus.counting_balanced_implies_zero}
Let \(I\) be a pure-dimensional incidence of dimension \(D\). On the counting
surface,
\[
\lambda(I)=0
\quad\Longleftrightarrow\quad
\text{\(I\) is closure-balanced in the sense of Axiom~\ref{ax:closure}}.
\]
Hence, on pure-dimensional incidences, exact counting closure is represented by
vanishing counting defect on that realized surface.
\end{proposition}

\begin{proof}
By Proposition~\ref{prop:pure-dimensional-counting-realization},
\(\lambda(I)=\binom{D}{2}-(D-1)!\). Therefore \(\lambda(I)=0\) exactly when
the pair count and order count agree, which is precisely Closure Balance in
the counting model used here.
\end{proof}

\begin{proposition}[Counting realization reflects terminality on the pure-dimensional surface]
\label{prop:counting-realization-reflects-terminality}\lean{CoherenceCalculus.counting_realization_reflects_terminality}
Let \(I\) be a pure-dimensional incidence of dimension \(D\). Assume the
counting realization is sound on this surface, in the sense that zero counting
defect represents zero bedrock closure defect. Then
\[
\lambda(I)=0
\quad\Longleftrightarrow\quad
\Lambda(I)=0
\quad\Longleftrightarrow\quad
I\text{ is terminal.}
\]
In particular, on pure-dimensional incidences, exact counting closure is the
realized surface of bedrock terminality.
\end{proposition}

\begin{proof}
By Proposition~\ref{prop:pure-dimensional-counting-realization},
\(\lambda(I)=0\) exactly when the counting closure residual vanishes. By the
soundness assumption on the pure-dimensional counting surface, this is
precisely the condition \(\Lambda(I)=0\). Terminality is then equivalent to
zero bedrock defect by Axiom~\ref{ax:termination}.
\end{proof}

\begin{center}
\begin{tabular}{cccc}
\toprule
$D$ & Pair Count & Order Count & $\cH(D)$ \\
\midrule
2 & 1 & 1 & 0 (Balanced) \\
3 & 3 & 2 & $+1$ (Residual) \\
4 & 6 & 6 & 0 (Balanced) \\
\bottomrule
\end{tabular}
\end{center}

\begin{theorem}[Uniqueness of the Residual Dimension]\label{thm:barrier-uniqueness}\lean{barrier_D3}
$D = 3$ is the unique dimension with $\cH(D) > 0$.
Equivalently, $D=3$ is the unique residual dimension.
\end{theorem}

(Proof in Appendix~\ref{sec:appendix-topology}.)

\begin{corollary}[Minimal nonterminal counting carrier]
\label{cor:minimal-nonterminal-counting-carrier}\lean{CoherenceCalculus.minimal_nonterminal_counting_carrier}
Among the nondegenerate counting dimensions, \(D=3\) is the minimal carrier on
which the counting realization of the closure defect remains nonzero:
\[
\lambda(I)=\cH(3)=1.
\]
Thus \(D=3\) is the first nonterminal mediation layer between the unit seed and
the balanced carrier.
\end{corollary}

\begin{proof}
By Theorem~\ref{thm:barrier-uniqueness}, \(D=3\) is the unique dimension with
\(\cH(D)>0\). Since \(\cH(3)=1\), the claim follows.
\end{proof}

\begin{theorem}[Unique nondegenerate balanced dimension in the counting model]\label{thm:forced-stable-dimension}\lean{stable_dimension_classification}\lean{closureStableDimension_globally_forced}
The only closure-balanced dimensions are $D = 2$ and $D = 4$.
Among these, $D = 4$ is the unique nontrivial balanced dimension.
\end{theorem}

\begin{proof}
Exact closure occurs at $D = 2$ and $D = 4$.
For $D = 0,1,3$ the residual is nonzero by direct computation, and for every $D \geq 5$ the order count strictly dominates the pair count.
Hence no balanced dimension exists beyond $D = 4$, and the only balanced dimension beyond the trivial case $D = 2$ is $D = 4$.
\end{proof}

\begin{remark}[Role of the \(D=4\) theorem]
Theorem~\ref{thm:forced-stable-dimension} is a counting-classification theorem
on the Part~I surface. It identifies the canonical nondegenerate balanced
local carrier used in the later structural and physical layers. The physical
reading begins only when the explicit Part~II three-law hypothesis is imposed
on that Part~I carrier.
\end{remark}

\begin{corollary}[First nontrivial terminal counting carrier]
\label{cor:first-nontrivial-terminal-counting-carrier}\lean{CoherenceCalculus.first_nontrivial_terminal_counting_carrier}
The first nontrivial carrier on which the counting realization of the closure
defect vanishes is \(D=4\):
\[
\lambda(I)=\cH(4)=0.
\]
Accordingly, \(D=4\) is the first nontrivial terminal carrier on the counting
surface.
\end{corollary}

\begin{proof}
Theorem~\ref{thm:forced-stable-dimension} identifies \(D=4\) as the unique
nontrivial balanced dimension. By
Proposition~\ref{prop:pure-dimensional-counting-realization}, balanced
dimensions are exactly those on which the counting realization vanishes.
\end{proof}

% ------------------------------------------------------------

\subsection{The Proof Chain}
\label{subsec:proof-chain}
% ------------------------------------------------------------

From one primitive and three axioms:

\begin{center}
\begin{tabular}{ll}
\toprule
\textbf{Primitive:} & Constancy \\
\textbf{Definition:} & Primitive \\
\textbf{Axiom I:} & Existence \\
\textbf{Definition:} & Incidence \\
\textbf{Axiom II:} & Incidence \\
\textbf{Axiom III:} & Closure Balance \\
\midrule
$\Downarrow$ & \\
Distinction & (explicit in the primitive interface) \\
Finiteness & (explicit in the incidence-count interface) \\
Finite choice (constructive) & (from explicit finite enumerations) \\
$\bbN$ and arithmetic & (from the canonical finite counting object) \\
Completeness: $\binom{D}{2} = (D-1)!$ & (from Closure Balance) \\
$D = 2, 4$ balanced; $D = 3$ residual & (from Completeness) \\
Canonical nondegenerate balanced carrier at $D = 4$ & (from global closure-balance classification) \\
\bottomrule
\end{tabular}
\end{center}

The operator calculus rests on this counting classification only at the moment
one passes to the canonical realization. The structural layer of Part~I does
not assume these coordinates in advance.

\paragraph{Connection to Coherence Calculus.}
In the canonical realization of the next section, the horizon projection
$\CCPi{h}$ implements the pair-count/order-count duality of Closure Balance at
each resolution level. The residual $\cH(D)$ becomes, after projection, a
defect term recording failure of closure at finite horizon. In this sense the
operator calculus organizes closure-balance failure under observation.

% --- Intrinsic decomposition and channel structure ---

\subsection{Primitive seed and local type profiles}

The counting derivation also fixes the local finite carrier on which any later
typed branch classification must act. At the bedrock level, this is not yet a
catalog of physical laws. It is a constructive statement that the local
closure algebra has only one primitive seed, one residual interface, and one
nontrivial balanced carrier.

\begin{definition}[Balanced closure slot count]
\label{def:balanced-closure-slot-count}
\lean{balancedClosureSlotCount}
For a balanced dimension $D$, define its \emph{closure slot count} by
\[
N_{\mathrm{cl}}(D) := \binom{D}{2} = (D-1)!.
\]
This is the common pair/order count carried by a balanced local incidence.
\end{definition}

\begin{definition}[Local type profile]
\label{def:local-type-profile}
\lean{LocalTypeProfile}
Let $D$ be balanced. A \emph{local type profile} at dimension $D$ is a finite
decomposition of the closure slot count into positive channel sizes:
\[
N_{\mathrm{cl}}(D)=n_1+\cdots+n_r,
\qquad
n_i\in\bbN,\quad n_i\ge 1.
\]
Equivalently, a local type profile is a partition of $N_{\mathrm{cl}}(D)$.
It records only the finite irreducible channel sizes. No symmetry group,
continuum field interpretation, or physical label is assumed at this stage.
\end{definition}

\begin{theorem}[Unique unit seed and six-slot balanced carrier]
\label{thm:unit-seed-and-six-slot-carrier}
\lean{unit_seed_and_six_slot_carrier}
The only balanced dimensions are $D=2$ and $D=4$, and their closure slot
counts are
\[
N_{\mathrm{cl}}(2)=1,
\qquad
N_{\mathrm{cl}}(4)=6.
\]
Consequently:
\begin{enumerate}[label=(\roman*), itemsep=0.2em]
\item $D=2$ furnishes the unique unit-balanced seed;
\item $D=4$ furnishes the unique nontrivial balanced local carrier, with six
      closure slots.
\end{enumerate}
\end{theorem}

(Proof in Appendix~\ref{sec:appendix-topology}.)

\begin{corollary}[Irreducibility of the primitive seed]
\label{cor:primitive-seed-irreducible}
\lean{primitive_seed_irreducible}
The only local type profile at the unit-balanced seed is
\[
1.
\]
In particular, the primitive seed admits no nontrivial decomposition into
smaller positive closure channels.
\end{corollary}

\begin{proof}
By Theorem~\ref{thm:unit-seed-and-six-slot-carrier},
$N_{\mathrm{cl}}(2)=1$. The only partition of $1$ into positive integers is
$1$ itself.
\end{proof}

\begin{corollary}[Single-channel unit seed]
\label{cor:single-channel-unit-seed}\lean{CoherenceCalculus.single_channel_unit_seed}
Any local law carried by the unit-balanced seed $D=2$ is necessarily
single-channel. In particular, no nontrivial local carrier splitting can occur
at the seed.
\end{corollary}

\begin{proof}
Corollary~\ref{cor:primitive-seed-irreducible} says that the only local type
profile at the seed is the one-part profile \(1\). Therefore every local law
carried by that seed has exactly one irreducible local channel.
\end{proof}

\begin{corollary}[Repeated transport at the seed introduces no new carrier data]
\label{cor:d2-repeat-no-new-carrier}\lean{CoherenceCalculus.repeat_cycle_monodromy_factorization}\lean{CoherenceCalculus.single_channel_unit_seed}
Let the unit-balanced seed \(D=2\) carry a reversible transport cycle \(W\) in
the sense of Definition~\ref{def:transport-cycle}. If a second cycle is a
repeat \(W^{\times m}\) with \(m>1\), then:
\begin{enumerate}[label=(\roman*), itemsep=0.2em]
\item its monodromy is \(\mathcal M(W)^m\);
\item it remains single-channel on the seed;
\item it creates no new local type profile and no new local carrier splitting.
\end{enumerate}
\end{corollary}

\begin{proof}
Part (i) is Proposition~\ref{prop:repeat-cycle-monodromy-factorization}. By
Corollary~\ref{cor:single-channel-unit-seed}, every local law on the seed is
single-channel, so the repeated cycle remains on that same one-channel carrier,
proving (ii). Since the seed admits no local type profile other than \(1\), no
repeat can create a new local type profile or a new carrier splitting, which is
(iii).
\end{proof}

\begin{corollary}[Primitive supports control all seed spectral data]
\label{cor:d2-primitive-support-profile}\lean{CoherenceCalculus.repeat_cycle_monodromy_factorization}\lean{CoherenceCalculus.PrimitiveTransportCycle}\lean{CoherenceCalculus.single_channel_unit_seed}
Let \(\mathcal C\) be any finite family of reversible transport cycles carried
by the unit-balanced seed \(D=2\). Then there exists a finite family
\[
\mathcal P
\]
of primitive transport cycles on that seed such that for every
\[
C\in\mathcal C
\]
there are \(P\in\mathcal P\) and \(m\ge 1\) with
\[
\mathcal M(C)=\mathcal M(P)^m.
\]
Consequently, every seed-side spectral datum is generated by primitive
supports together with repeat exponents, and no non-primitive cycle introduces
a new seed carrier type.
\end{corollary}

(Proof in Appendix~\ref{sec:appendix-topology}.)

\begin{theorem}[Residual mediation and finite local type profiles]
\label{thm:finite-local-type-profiles}
\lean{finite_local_type_profiles}
The passage from the unit-balanced seed to the nontrivial balanced carrier is
mediated by the unique residual dimension $D=3$. Moreover every local type
profile on the canonical balanced carrier $D=4$ is one of the finitely many
partitions of $6$:
\[
\begin{gathered}
6,\quad
5+1,\quad
4+2,\quad
4+1+1,\quad
3+3,\quad
3+2+1,\\
3+1+1+1,\quad
2+2+2,\quad
2+2+1+1,\quad
2+1+1+1+1,\quad
1+1+1+1+1+1.
\end{gathered}
\]
Consequently, bedrock admits only finitely many local type profiles before any
additional symmetry, admissibility, or realization hypotheses are imposed.
\end{theorem}

(Proof in Appendix~\ref{sec:appendix-topology}.)

\begin{corollary}[Choice-free reduction of local type classification]
\label{cor:choice-free-type-classification}
\lean{choice_free_type_classification}
Any later typed branch classification in this calculus reduces first to
the finite counting data
\[
\text{unit seed }(1),\qquad
\text{residual interface }(D=3),\qquad
\text{balanced carrier }(6),
\]
and then to a choice among the finitely many local type profiles of
Theorem~\ref{thm:finite-local-type-profiles}. No open-ended continuum of
primitive types is introduced at bedrock level.
\end{corollary}

\begin{proof}
The seed/interface/carrier trichotomy is exactly the content of
Theorem~\ref{thm:unit-seed-and-six-slot-carrier} together with
Theorem~\ref{thm:finite-local-type-profiles}. The final sentence records that
the possible local profiles have already been explicitly enumerated.
\end{proof}

\begin{corollary}[Residual localization on the unique interface]
\label{cor:residual-localization-on-d3}
\lean{barrier_D3, finite_local_type_profiles}
Any later local realization on the canonical balanced carrier that singles out
one residual mediation layer compatible with the counting bedrock must localize
that layer on the unique residual interface \(D=3\). There is no second
counting-compatible local residual interface between the unit seed and the
nontrivial balanced carrier.
\end{corollary}

\begin{proof}
By Theorem~\ref{thm:finite-local-type-profiles}, the passage from the
unit-balanced seed to the nontrivial balanced carrier is mediated by the unique
residual dimension \(D=3\). Therefore any local realized extension that marks a
single residual mediation layer while remaining compatible with the counting
bedrock must mark that same interface. No second such interface exists in the
seed/interface/carrier trichotomy of
Corollary~\ref{cor:choice-free-type-classification}.
\end{proof}

\begin{definition}[Logical type signature]
\label{def:logical-type-signature}
\lean{LogicalTypeSignature}
A \emph{logical type signature} on the canonical balanced carrier is a local
type profile at $D=4$. Equivalently, it is a partition of the six-slot closure
count
\[
6=n_1+\cdots+n_r,
\qquad
n_i\in\bbN,\quad n_i\ge 1.
\]
The parts $n_i$ record only irreducible local channel sizes.
\end{definition}

\begin{definition}[Symmetry-refined logical type]
\label{def:symmetry-refined-logical-type}
\lean{SymmetryRefinedLogicalType}
Let $V$ be a six-dimensional local carrier equipped with a finite unitary
symmetry action of a group $G$, and decompose $V$ into its nonzero isotypic
components:
\[
V=\bigoplus_{\rho}\,V_\rho.
\]
The \emph{symmetry-refined logical type} of $(V,G)$ is the multiset of positive
dimensions
\[
\{\dim V_\rho : V_\rho\neq 0\}.
\]
\end{definition}

\begin{proposition}[Every symmetry-refined logical type lies over a finite logical signature]
\label{prop:symmetry-refined-logical-type}
\lean{symmetry_refined_logical_type}
Let $(V,G)$ be as in Definition~\ref{def:symmetry-refined-logical-type}. Then:
\begin{enumerate}[label=(\roman*), itemsep=0.2em]
\item the symmetry-refined logical type is a partition of $6$;
\item its number of distinct symmetry channels is at most $6$;
\item its underlying logical type signature is one of the finitely many
      signatures listed in Theorem~\ref{thm:finite-local-type-profiles}.
\end{enumerate}
\end{proposition}

(Proof in Appendix~\ref{sec:appendix-topology}.)

\begin{corollary}[Possible logical types are finite at bedrock level]
\label{cor:logical-types-finite}
\lean{logical_types_finite}
At bedrock level, possible logical types are
finite and constructive: they are exactly the logical type signatures of
Definition~\ref{def:logical-type-signature}, together with their
symmetry-refined realizations of
Definition~\ref{def:symmetry-refined-logical-type}. In particular, no later
typed branch can introduce more than six irreducible local channels on the
canonical balanced carrier.
\end{corollary}

\begin{proof}
The finite list of logical type signatures is given by
Theorem~\ref{thm:finite-local-type-profiles}. Any symmetry-refined realization
of such a type lies over one of those signatures by
Proposition~\ref{prop:symmetry-refined-logical-type}. The final sentence is
exactly Proposition~\ref{prop:symmetry-refined-logical-type}(ii).
\end{proof}

\begin{definition}[Local slot carrier of a balanced four-incidence]
\label{def:local-slot-carrier}
\lean{LocalSlot}\lean{LocalSlotCarrier}
Fix a balanced incidence with primitive set $\{1,2,3,4\}$. Its \emph{local slot
set} is the set of unordered primitive pairs
\[
\mathcal S_4:=\{\{i,j\}:1\le i<j\le 4\},
\]
so $|\mathcal S_4|=6$. The associated \emph{local slot carrier} is the
six-dimensional real vector space
\[
V_{\mathrm{slot}}:=\bbR[\mathcal S_4]
\]
with basis vectors $e_{ij}=e_{\{i,j\}}$ indexed by the six slots.
\end{definition}

\begin{definition}[Intrinsic relabeling action]
\label{def:intrinsic-relabeling-action}
\lean{Primitive4}\lean{IntrinsicRelabelingAction}
The full relabeling group of the four primitives is the symmetric group
$S_4$. It acts on the slot set by
\[
\sigma\cdot \{i,j\}:=\{\sigma(i),\sigma(j)\},
\qquad
\sigma\in S_4,
\]
and therefore linearly on $V_{\mathrm{slot}}$ by
\[
\sigma\cdot e_{ij}:=e_{\sigma(i)\sigma(j)}.
\]
\end{definition}

\paragraph{No new axiom enters here.}
\label{rem:no-new-axiom-s4}
Definition~\ref{def:intrinsic-relabeling-action} introduces no additional
bedrock axiom. Once a four-incidence is fixed, the set of all bijections of
its four primitives is canonically $S_4$. The only extra hypothesis that could
enter later is a realization-level equivariance assumption asserting that a
particular operator law commutes with this intrinsic relabeling action.

\begin{theorem}[Canonical intrinsic decomposition of the six-slot carrier]
\label{thm:intrinsic-six-slot-decomposition}
\lean{intrinsic_six_slot_decomposition}
Under the intrinsic relabeling action of $S_4$, the six-slot carrier admits a
canonical direct-sum decomposition
\[
V_{\mathrm{slot}}=U_1\oplus U_2\oplus U_3
\]
with
\[
\dim U_1=1,\qquad
\dim U_2=2,\qquad
\dim U_3=3.
\]
Consequently the intrinsic symmetry-refined logical type of the canonical
balanced carrier is forced to be
\[
1+2+3.
\]
\end{theorem}

(Proof in Appendix~\ref{sec:appendix-topology}.)

\begin{corollary}[Intrinsic symmetry removes the remaining signature ambiguity]
\label{cor:intrinsic-symmetry-forces-123}
\lean{intrinsic_symmetry_forces_123}
If logical types are classified by the full intrinsic relabeling symmetry of a
balanced four-incidence, then the finite list of bedrock signatures collapses
to a unique canonical signature:
\[
1+2+3.
\]
\end{corollary}

\begin{proof}
Theorem~\ref{thm:intrinsic-six-slot-decomposition} constructs the unique
intrinsic $S_4$-invariant channel decomposition of the six-slot carrier. Its
dimensions are exactly $1$, $2$, and $3$.
\end{proof}

\begin{theorem}[Algebraic closure of the edge carrier under exterior product]
\label{thm:fermionic-sector-from-ext2}\lean{CoherenceCalculus.ExteriorSquareFermionicData, CoherenceCalculus.fermionic_sector_from_ext2}
Let $V_{\mathrm{slot}} = U_1 \oplus U_2 \oplus U_3$ be the canonical
$S_4$-decomposition of the edge carrier
(Theorem~\ref{thm:intrinsic-six-slot-decomposition}). Then:
\begin{enumerate}[label=\textup{(\alph*)}, itemsep=0.3em]
\item $\mathrm{Ext}^2(U_2) \cong \mathrm{sign}$, the sign representation
  of~$S_4$ (dimension~$1$).
\item $\mathrm{Ext}^2(U_3) \cong \mathrm{sign} \otimes \mathrm{std}$, the
  sign-twisted standard representation of~$S_4$ (dimension~$3$).
\item The edge carrier $V_{\mathrm{slot}}$ contains only the irreducible
  representations $\mathrm{trivial}(1) \oplus \mathrm{2\text{-}dim}(2)
  \oplus \mathrm{std}(3)$. The $\mathrm{sign}(1)$ and
  $\mathrm{sign} \otimes \mathrm{std}(3)$ representations are absent.
\item The antisymmetric products of the bosonic channels therefore generate
  representations not contained in the edge carrier. The algebraic closure
  of $V_{\mathrm{slot}}$ under $\mathrm{Ext}^2$ is the full $S_4$
  representation ring restricted to these five irreducibles, of total
  dimension $1 + 2 + 3 + 1 + 3 = 10$.
\end{enumerate}
The additional $4$ dimensions ($\mathrm{sign} \oplus
\mathrm{sign} \otimes \mathrm{std}$) carry the sign character of~$S_4$:
under any odd permutation, they acquire a factor of~$-1$.
\end{theorem}

\begin{proof}
For part~(a): $U_2$ is the $2$-dimensional irreducible representation
of~$S_4$. Its exterior square $\mathrm{Ext}^2(U_2)$ has dimension
$\binom{2}{2} = 1$. The character of $\mathrm{Ext}^2(U_2)$ at a group
element~$g$ is $\frac{1}{2}(\chi_{U_2}(g)^2 - \chi_{U_2}(g^2))$. Direct
computation from the $S_4$ character table gives
$\chi_{\mathrm{Ext}^2(U_2)} = \chi_{\mathrm{sign}}$.

For part~(b): $U_3 = \mathrm{std}$ has dimension~$3$.
$\mathrm{Ext}^2(U_3)$ has dimension $\binom{3}{2} = 3$. The same character
formula gives
$\chi_{\mathrm{Ext}^2(\mathrm{std})} = \chi_{\mathrm{sign} \otimes
\mathrm{std}}$.

Parts~(c) and~(d) follow by comparing the irreducible content of
$V_{\mathrm{slot}}$ with the outputs of the $\mathrm{Ext}^2$ operations.
The edge carrier contains trivial, $2$-dim, and std. The $\mathrm{Ext}^2$
products generate sign and $\mathrm{sign}\otimes\mathrm{std}$, which are
absent from the edge carrier.
The total dimension of the five irreducibles is
$1 + 1 + 2 + 3 + 3 = 10$.
\end{proof}

\begin{remark}[Observation: the sign character and fermions]
\label{rem:fermionic-observation}\lean{CoherenceCalculus.ExteriorSquareFermionicData, CoherenceCalculus.fermionic_sector_from_ext2}
The additional representations $\mathrm{sign}(1)$ and
$\mathrm{sign} \otimes \mathrm{std}(3)$ both carry the sign character
of~$S_4$: they change sign under odd permutations. Under the continuum
identification, permutations become spatial coordinate exchanges, and the
sign character becomes the Pauli exclusion signature. The $1 + 3$ split of
the fermionic sector is suggestive of the lepton/quark decomposition
observed in the Standard Model~\cite{PeskinSchroeder1995}. This observation is recorded here as
structural; the physical identification is not claimed as a theorem of the
present manuscript.
\end{remark}

\begin{theorem}[The six-slot carrier is the symmetric square of the primitive \(3\)-channel]
\label{thm:six-slot-carrier-is-symmetric-square-of-u3}
\lean{six_slot_carrier_is_symmetric_square_of_u3}
Let \(U_3\subseteq V_{\mathrm{slot}}\) be the primitive three-dimensional
channel of Theorem~\ref{thm:intrinsic-six-slot-decomposition}. Then there is a
canonical symmetric bilinear map
\[
\beta_{\mathrm{slot}}:U_3\times U_3\to V_{\mathrm{slot}}
\]
whose image generates \(V_{\mathrm{slot}}\). Consequently the induced linear map
\[
\overline\beta_{\mathrm{slot}}:\mathrm{Sym}^2(U_3)\to V_{\mathrm{slot}}
\]
is an isomorphism. In particular,
\[
V_{\mathrm{slot}}\cong \mathrm{Sym}^2(U_3)
\]
canonically, compatibly with the intrinsic relabeling action.
\end{theorem}

(Proof in Appendix~\ref{sec:appendix-topology}.)

\begin{definition}[Intrinsic adjacency classes on slot pairs]
\label{def:intrinsic-slot-adjacency}
\lean{IntrinsicSlotAdjacency}\lean{intrinsicSlotAdjacency}
Let $\alpha,\beta\in\mathcal S_4$ be slots of a balanced four-incidence.
They are:
\begin{enumerate}[label=(\roman*), itemsep=0.2em]
\item \emph{equal} if $\alpha=\beta$;
\item \emph{adjacent} if $\alpha\neq\beta$ and $\alpha\cap\beta$ contains one
      primitive;
\item \emph{disjoint} if $\alpha\cap\beta=\varnothing$.
\end{enumerate}
These are exactly the three orbit types of ordered slot pairs under the
intrinsic relabeling action of $S_4$.
\end{definition}

\begin{theorem}[Intrinsic normal form of the six-slot operator algebra]
\label{thm:intrinsic-slot-normal-form}
\lean{intrinsic_slot_normal_form}
Let $A:V_{\mathrm{slot}}\to V_{\mathrm{slot}}$ be a linear operator commuting
with the intrinsic relabeling action of $S_4$. Then there exist unique scalars
$a,b,c\in\bbR$ such that for every slot $\alpha\in\mathcal S_4$,
\[
A e_\alpha
=
a\,e_\alpha
+
b\sum_{\beta\sim\alpha} e_\beta
+
c\sum_{\beta\perp\alpha} e_\beta,
\]
where $\beta\sim\alpha$ denotes the four slots adjacent to $\alpha$ and
$\beta\perp\alpha$ denotes the unique slot disjoint from $\alpha$.
Equivalently, in the slot basis, the matrix entries of $A$ depend only on the
three intrinsic classes of Definition~\ref{def:intrinsic-slot-adjacency}:
\[
A_{\alpha\beta}=
\begin{cases}
a,& \alpha=\beta,\\
b,& \alpha,\beta\text{ adjacent},\\
c,& \alpha,\beta\text{ disjoint}.
\end{cases}
\]
\end{theorem}

(Proof in Appendix~\ref{sec:appendix-topology}.)

\begin{corollary}[Intrinsic channel eigenvalues]
\label{cor:intrinsic-channel-eigenvalues}
\lean{intrinsic_channel_eigenvalues}
Let $A$ be as in Theorem~\ref{thm:intrinsic-slot-normal-form}, with parameters
$a,b,c$. Then the canonical channels of
Theorem~\ref{thm:intrinsic-six-slot-decomposition} are $A$-invariant, and $A$
acts on them by the scalar eigenvalues
\[
\lambda_1=a+4b+c,\qquad
\lambda_2=a-2b+c,\qquad
\lambda_3=a-c
\]
on the $1$-, $2$-, and $3$-dimensional channels respectively.
\end{corollary}

(Proof in Appendix~\ref{sec:appendix-topology}.)

\begin{corollary}[Intrinsic-equivariant laws are three-parameter laws]
\label{cor:intrinsic-equivariant-three-parameter}
\lean{intrinsic_equivariant_three_parameter}
Any local law on the canonical balanced carrier that commutes with the full
intrinsic relabeling action of the four-incidence is determined by exactly
three scalar channel parameters. Equivalently, such a law preserves the forced
$1+2+3$ logical-type decomposition and cannot mix those three channels.
\end{corollary}

\begin{proof}
Theorem~\ref{thm:intrinsic-slot-normal-form} shows that an intrinsic-equivariant
local law is determined by the three scalars $a,b,c$. Corollary
~\ref{cor:intrinsic-channel-eigenvalues} shows that the corresponding operator
acts separately on the forced $1$-, $2$-, and $3$-dimensional channels.
\end{proof}

\paragraph{Relation to symmetry reduction.}
\label{rem:type-profile-symmetry-reduction}
Once a concrete finite symmetry action is imposed on the canonical balanced
carrier, the symmetry reduction theorems of Chapter~4 refine one of the local
type profiles above into isotypic operator blocks and multiplicity spaces.
The counting bedrock therefore fixes the finite carrier size first; symmetry
and admissibility only decide how that already forced carrier decomposes.
